\documentclass[a4paper,11pt]{article}
\usepackage[utf8]{inputenc}
\usepackage{amsmath}
\usepackage{amssymb}
\usepackage{amsthm}
\usepackage{amsfonts}
\usepackage{amsbsy}
\usepackage{esint}
\usepackage{mathtools}
\usepackage{array}
\usepackage{tikz}
\usepackage{url}
\usepackage{graphicx}
\usepackage{physics}
\usepackage{dsfont}
\usepackage{xcolor}
\usepackage[colorlinks=true, allcolors=blue]{hyperref}
\newcommand{\Schatten}{{\mathfrak S}}
\newcommand{\per}{\mathrm{per}}
\newcommand{\Ran}{\mathrm{Ran}}
\newcommand{\Trper}{\underline{\Tr}}
\newcommand{\BZ}{\mathcal{B}}
\theoremstyle{plain}
\newtheorem{theorem}{Theorem}[section]
\newtheorem{definition}[theorem]{Definition}
\newtheorem{assumption}[theorem]{Assumption}
\newtheorem{lemma}[theorem]{Lemma}
\newtheorem{corollary}[theorem]{Corollary}
\newtheorem{proposition}[theorem]{Proposition}
\newtheorem{hyp}[theorem]{Assumption}
\theoremstyle{remark}
\newtheorem{remark}[theorem]{Remark}
\newtheorem*{notation}{Notation}
\newtheorem{example}[theorem]{Example}
\renewcommand{\vec}[1]{#1}
\usepackage{tikz-cd}

\makeatletter
\tikzset{
  column sep/.code=\def\pgfmatrixcolumnsep{\pgf@matrix@xscale*(#1)},
  row sep/.code   =\def\pgfmatrixrowsep{\pgf@matrix@yscale*(#1)},
  matrix xscale/.code=
    \pgfmathsetmacro\pgf@matrix@xscale{\pgf@matrix@xscale*(#1)},
  matrix yscale/.code=
    \pgfmathsetmacro\pgf@matrix@yscale{\pgf@matrix@yscale*(#1)},
  matrix scale/.style={/tikz/matrix xscale={#1},/tikz/matrix yscale={#1}}}
\def\pgf@matrix@xscale{1}
\def\pgf@matrix@yscale{1}
\makeatother


\usepackage[margin=2.0cm]{geometry}


\def\R{{\mathbb R}}
\def\N{{\mathbb N}}
\def\C{{\mathbb C}}
\def\Z{{\mathbb Z}}
\def\eps{{\varepsilon}}

\def\cB{{\mathcal B}}
\def\cC{{\mathcal C}}
\def\cH{{\mathcal H}}
\def\cL{{\mathcal L}}
\def\cN{{\mathcal N}}
\def\cR{{\mathcal R}}


\def\1{{\mathds 1}}

\def\dps{\displaystyle }


\newcommand{\sub}[1]{_{\mathrm{#1}}}
\newcommand{\su}[1]{^{\mathrm{#1}}}


\title{Conductivity at conical crossings}
\author{}
\date{}
\begin{document}
\maketitle

\begin{abstract}
We consider independent electrons in a periodic crystal in their ground
state, and turn on a uniform electric field at some prescribed time.
We rigorously define the current per unit area and study the static
conductivity of a two-dimensional semimetal. The conductivity tensor is isotropic and each
conical intersection contributes as $\frac{1}{16}$ to the total conductivity.
\end{abstract}
\tableofcontents
\section{Introduction}

We consider a two-dimensional perfect crystal with periodic lattice
$\cR\subset\R^2$. We use atomic units and a spinless model, with Hamiltonian
\begin{equation}\label{eq:periodic_Hamiltonian}
H=-\frac12\Delta+V\quad\hbox{on }L^2(\R^2;\C),
\end{equation}
where $V\in H^1_{\rm per}(\Omega;\R)$ and $\Omega$ is a unit cell of $\cR$.
At zero temperature, the ground-state density matrix is given by
\begin{equation}\label{eq:GSDM}
\gamma_0=\mathds 1(H\le\mu_{\rm F}),
\end{equation}
where $\mu_{\rm F}\in\R$ is the Fermi level.

The purpose of this article is to analyze mathematically the behavior
of the current appearing in the crystal when a uniform external field
is turned on instantaneously at the initial time $t=0$.
Fix $e_\beta,e_\alpha\in\R^2$ once and for all, for the field and observation
directions respectively. They need not be normalized, linearly independent,
or orthogonal. Put $x_\beta=x\cdot e_\beta$ and $p=-i\nabla$. For $\eps\ge0$,
\begin{align}\label{def:Heps}
H_\eps=H-\eps x_\beta,
\end{align}
where $H_\eps$ denotes the self-adjoint closure of its restriction to
$C_c^\infty(\R^2;\C)$ (see Proposition~\ref{prop:setting}). The state at
$t\ge0$ is
\begin{equation}\label{def:gammaeps}
\gamma_\eps(t)=e^{-itH_\eps}\gamma_0e^{itH_\eps}.
\end{equation}
The current is defined as
\begin{equation}\label{eq:current}
j_\eps(t)=\Trper(J\gamma_\eps(t)),
\end{equation}
where $\Trper$ is the trace per unit area, defined in
Section~\ref{sec:notation}, and
\begin{equation}\label{def:J}
J=p\cdot e_\alpha.
\end{equation}


For time-reversible 2D semimetals such as graphene, the time-averaged conductivity
\begin{align}\label{def:conductivity}
\sigma = \lim_{t\to\infty}\frac1t\int_0^t
\lim_{\eps\to0^+}\frac{j_\eps(t')}{\eps}\,dt'
\end{align}
has a finite value equal to $\frac1{16}e_\alpha\cdot e_\beta$ times
the number of Dirac points in the Brillouin zone.


Note that our results use an averaging in time, and we are unable to
conclude anything about what would be the naive definition of the
conductivity
\begin{align}
  \label{eq:conductivity_naive}
  \lim_{t \to  \infty} \lim_{\eps \to 0^+} \frac{j_\eps(t)}{\eps}.
\end{align}
A form of averaging of time fluctuations is always necessary to infer
zero-frequency behavior from step responses in non-dissipative
systems, even in the linear case. The easiest way to see this is by
the very simple model for the relationship between an input $I(t)$ and
an output $O(t)$:
  \begin{equation} \label{eq:toymodel}
    i \dot O(t) = \omega O(t) + I(t).
  \end{equation}
  This simplified model describes a forced oscillator with eigenfrequency $\omega$, and arises from the linear response of the time-dependent Schr\"odinger equation of a two-level system. For a
  constant input~$I_{0}$, there is a steady state solution
  $O_0 = \widehat R_0 I_{0}$, where $\widehat R_0 = -\frac 1 \omega$ is
  the zero-frequency transfer function of the system. However, since
  this system is oscillatory, this steady state may never be reached:
  if $I$ is brutally switched on at time $0$ with $I(t) = O(t) = 0,$
  if $t \leq 0$, $I(t) = I_{0}$ if $t > 0$, then
  $ O(t) = O_0(1- {\rm e}^{-i\omega t}) = \widehat{R}_0 I_{0} (1- {\rm
    e}^{-i\omega t})$ and we cannot define $\widehat R_0$ as the limit
  of $O(t)/I_{0}$ when $t$ goes to infinity. However, by averaging we
  obtain
  \begin{align*}
     \lim_{t\to\infty}\frac 1 t  \int_{0}^{t} \frac{O(t')}{I_{0}} \, dt' = \widehat R_0.
   \end{align*}

In our results in the linear response regime, we consider the quantity
\begin{align*}
  j^{\rm LR}(t) = \lim_{\eps \to 0^+} \frac {j_\eps(t)}{\eps}.
\end{align*}
and then compute the conductivity in the limit $t \to \infty$. This
order of limits is necessary to ensure that the electronic state never
deviates significantly from its equilibrium. Considering the opposite
limit, i.e. the infinite-time dynamics of
$j_\eps(t)$ at $\eps$ fixed, is an extremely
hard problem, as it is for any dynamics of non-dissipative systems.
Our method of proof is based on the standard gauge change
$\widetilde\psi(x,t)=e^{-i\eps t x_\beta}\psi(x,t)$
that transforms the constant in time but non-spatially-periodic
Hamiltonian $H_\eps$ into the time-dependent Hamiltonian
$\widetilde H_\eps(t)=\frac12(p+\eps t e_\beta)^2+V$.
This Hamiltonian is spatially periodic, and the study of its dynamics
can be reduced via Bloch--Floquet theory to that of its fibers
$H_{k+\eps t e_\beta}$ acting on periodic functions
(Section~\ref{sec:preliminaries}). Fiber by fiber, this time-dependent
Hamiltonian can then be treated using the tools of time-dependent
perturbation theory (Section~\ref{sec:linear-response}). For each fixed $t$,
we use linear response to expand $j_\eps(t)$ to first order in $\eps$
(Proposition~\ref{prop:response}).


The structure of the paper is as follows. We describe our results in
Section~\ref{sec:results}: we define the current in
Proposition~\ref{prop:setting}, and state Theorem~\ref{theo:semi-metal}.
We devote Section~\ref{sec:preliminaries} to preliminaries about the
regularity and Bloch decomposition of the current.
Section~\ref{sec:linear-response} states and proves the linear response result.
Section~\ref{sec:semimet} is devoted to the proof of the conductivity
formula.

\section{Main results: electrical current in periodic materials}\label{sec:results}


\subsection{Notation}
\label{sec:notation}
We denote by $\cR^*$ the dual lattice
\[
\cR^*=\{K\in\R^2:K\cdot R\in2\pi\Z\hbox{ for all }R\in\cR\},
\]
by $\Omega$ a unit cell in physical space, and by $\cB$ a unit cell in
reciprocal space, called the Brillouin zone. We write
$\partial_\alpha=e_\alpha\cdot\nabla_k$ and
$\partial_\beta=e_\beta\cdot\nabla_k$ for the two fixed directional derivatives.



The {\it $\cR$-periodic Lebesgue and
Sobolev spaces} are denoted by
\begin{align*}
L^p_{\rm per}&:=\left\{ u \in L^p_{\rm loc}(\R^2;\C) \; | \; u \mbox{ $\cR$-periodic} \right\}, \\
H_{\rm per}^{s}&:=\left\{ u \in H_{\rm loc}^{s}(\R^2;\C) \; | \; u \mbox{ $\cR$-periodic} \right\}.
\end{align*}
The space of bounded linear operators on a Hilbert space~$\mathcal{H}$ is denoted by ${\mathcal L}(\mathcal{H})$, and the {Schatten class} of bounded operators~$A \in {\mathcal L}(\mathcal{H})$ such that $\Tr(|A|^p) < \infty$ by $\Schatten_{p}(\mathcal{H})$.
For $R \in \cR$, we denote by~$\tau_R$ the {translation operator} formally defined by $\tau_R \phi = \phi(\cdot - R)$. Depending on the context, $\tau_R$ will be seen as a unitary operator on $L^2(\R^2;\C)$, or as a linear operator on some $\cR$-translation invariant subspace of ${\mathcal D}'(\R^2;\C)$.
A bounded operator on $L^2(\R^2;\C)$ is called $\cR$-periodic if it commutes with~$\tau_R$ for all $R \in \cR$. An unbounded self-adjoint operator on $L^2(\R^2;\C)$ is called $\cR$-periodic if its resolvent is $\cR$-periodic. A bounded $\cR$-periodic operator $A \in {\mathcal L}(L^2(\R^2;\C))$ is called locally trace-class if $\chi A \chi \in \Schatten_{1}(L^2(\R^2;\C))$ for any compactly supported function $\chi \in L^\infty(\R^2;\C)$. For $p \ge 1$, we denote by $\Schatten_{p,\per}$ the space of $\cR$-periodic operators $A \in {\mathcal L}(L^2(\R^2;\C))$ such that $|A|^p$ is locally trace class. Any operator $A \in \Schatten_{1,\per}$ has a density $\rho_A \in L^1_{\rm per}$ characterized by
$$
\forall \chi \in C^\infty_{\rm c}(\R^2;\C), \quad \Tr(A\chi) = \int_{\R^2} \rho_A \chi.
$$
The trace per unit area of an operator $A \in \Schatten_{1,\per}$ is defined as
$$
\Trper(A) = \frac{1}{|\Omega|} \Tr_{L^2(\R^2;\C)}(\1_\Omega A \1_\Omega) =
\fint_{\Omega}  \rho_A,
$$
where $\1_\Omega$ is the characteristic function of the unit cell
$\Omega$, and $\fint_\Omega$ is a shorthand notation for $\frac 1
{|\Omega|} \int_{\Omega}$. This formula is independent of the choice
of the unit cell $\Omega$.


\medskip

Since we are dealing here with periodic materials, we will use the {Bloch transform} (also
called Bloch-Floquet transform) \cite{ReedSimon4}. For $K \in \cR^\ast$, let $T_K$ be the unitary multiplication operator on $L^2_{\rm per}$ defined by
$$
\forall v \in L^2_{\rm per}, \quad (T_Kv)(x) = {\rm e}^{-iK\cdot x} v(x) \quad \mbox{ for a.a. } x \in \R^2,
$$
and
\begin{align*}
&L^2_{\rm qp}(L^2_{\rm per})\\
&:=\left\{ \R^2 \ni k \mapsto u_k \in L^2_{\rm per} \; | \; \int_\cB \|u_k\|_{L^2_{\rm per}}^2 \, dk < \infty, \; u_{k+K}=T_K u_k \mbox{ for all } K \in \cR^\ast \mbox{ and a.a. } k \in \R^2 \right\},
\end{align*}
the Hilbert space of $\cR^\ast$-quasi-periodic $L^2_{\rm per}$-valued functions on $\R^2$ endowed with the inner product
$$
\langle u,v\rangle_{L^2_{\rm qp}(L^2_{\rm per})}= \fint_\cB \langle u_k,v_k\rangle_{L^2_{\rm per}} \, dk.
$$
Here and below, the subscript qp refers to the quasi-periodicity property.
The Bloch transform then is the unitary map from $L^2(\R^2;\C)$ to $L^2_{\rm qp}(L^2_{\rm per})$ defined for $u \in C^\infty_{\rm c}(\R^2;\C)$ by
\begin{equation}\label{eq:Bloch_transform}
\forall k \in \R^2, \quad \forall x \in \R^2, \quad u_k(x) = \sum_{R \in \cR} u(x+R) \, {\rm e}^{-ik \cdot (x+R)}.
\end{equation}
Its inverse is given by
\begin{equation}\label{eq:inverse_Bloch_transform}
u(x) =\fint_\cB {\rm e}^{ik\cdot x} u_{k}(x) \, dk, \quad \mbox{for a.a. } x \in \R^2.
\end{equation}
Any $\cR$-periodic
operator $A \in {\mathcal L}(L^2(\R^2;\C))$ is decomposed by the Bloch
transform in the sense that there exists a function
$k \mapsto A_k$ in $L^\infty_{\rm qp}({\mathcal L}(L^2_{\rm per}))$ such that for any
$u \in L^2(\R^2;\C)$ and almost all $k \in \R^2$, $(Au)_k = A_ku_k$, and
\begin{equation}\label{eq:qpBloch}
A_{k+K} = T_K A_k T_K^\ast, \quad \mbox{ for all } K \in \cR^\ast \mbox{ and a.a. } k \in \R^2.
\end{equation}
The $A_{k}$'s are called the fibers of the operator $A$.
If $A \in \Schatten_{1,\per}$, then the function $k \mapsto A_{k}$ is
in $L^1_{\rm qp}(\Schatten_{1}(L^2_{\rm per}))$, the function $k
\mapsto \Tr(A_k)$ is in $L^1_{\rm loc}(\R^2)$, $\cR^{*}$-periodic, and we have
$$
\Trper(A) =(2\pi)^{-2} \int_\cB \Tr(A_k) \, dk.
$$
The Bloch decomposition theorem can be extended to unbounded
$\cR$-periodic self-adjoint operators using the resolvent \cite{ReedSimon4}.

\medskip


In the case of the periodic Hamiltonian operator $H$ given by \eqref{eq:periodic_Hamiltonian}, we have
  \begin{align} \label{eq:Hk}
    H_k = \frac 1 2 (-i\nabla+k)^2 + V.
  \end{align}
  For each $k \in \R^2$, $H_k$ is a bounded below self-adjoint operator on $L^{2}_{\per}$ with domain $H_{\per}^{2}$ and compact resolvent.
  Let $(\lambda_{n,k})_{n \in \N^\ast}$ be the non-decreasing sequence of eigenvalues of $H_k$ counting multiplicities
  $$
  \lambda_{1,k}\leq \lambda_{2,k}\leq \lambda_{3,k} \leq \cdots , \quad \lim_{n \to \infty} \lambda_{n,k}=+\infty,
  $$
  and we use the convention $\lambda_{0,k} = -\infty$. We denote by
  $(u_{n,k})_{n \in \N^\ast} \in (H_{\per}^{2})^{\mathbb N^{*}}$ an $L^2_{\rm per}$-orthonormal basis of associated eigenfunctions:
  \begin{align*}
    H_k u_{n,k} = \lambda_{n,k} u_{n,k}, \quad \langle u_{m,k},u_{n,k}\rangle_{L^2_{\rm per}}=\delta_{m,n}.
  \end{align*}
For $n\ge1$ and $k\in\R^2$, set
\[
P_{n,k}=\mathds 1(H_k\le\lambda_{n,k}),\qquad P_{0,k}=0.
\]
Whenever $\lambda_{n,k}<\lambda_{n+1,k}$, $P_{n,k}$ is the spectral
projector associated with the lowest $n$ eigenvalues, counting multiplicities:
\begin{equation}\label{eq:def_PNk}
P_{n,k}=\sum_{m=1}^n|u_{m,k}\rangle\langle u_{m,k}|.
\end{equation}

  Since $H_k$ is quasi-periodic, so is $P_{n,k}$, and the eigenvalues $\lambda_{n,k}$ are
  ${\mathcal R}^\ast$-periodic functions of $k$. By a min-max argument
  (see e.g. \cite{ReedSimon4}), there exists
  $\underline{C_1}, \overline{C_1} \in \R$, and
  $\underline{C_2},\overline{C_2}>0$ such that
  \begin{align}
    \label{eq:asymptotics_eigenvalues}
 \underline{C_1}+ \underline{C_2} n  \leq \lambda_{n,k} \leq  \overline{C_1}+ \overline{C_2} n.
\end{align}
The ground-state density matrix $\gamma_0$ has Bloch fibers
\begin{equation}\label{eq:gk0}
P_k=\mathds 1(H_k\le\mu_{\rm F}).
\end{equation}
Their ranks are uniformly bounded by the eigenvalue estimate above.
The current operator $J$ is also $\cR$-periodic, with fibers
\[
J_k=(p+k)\cdot e_\alpha=\partial_\alpha H_k.
\]


Lastly, for each $q \in \R^2$, we denote by $G_q$ the unitary multiplication operator on $L^2(\R^2;\C)$ defined by
\begin{equation}\label{eq:defGq}
\forall u \in L^2(\R^2;\C), \quad (G_qu)(x) = {\rm e}^{iq\cdot x} u(x) \quad \mbox{ for a.a. } x \in \R^2.
\end{equation}
The operator $G_q$ is not $\cR$-periodic, except
when $q \in \cR^\ast$ (in which case~$G_{q}$ is fibered, with
$G_{q,k} = T_{-q}$ for all $k$). However, for any $\cR$-periodic
operator $A \in {\mathcal L}(L^2(\R^2;\C))$ and any $q \in \R^2$,
the operator $G_qAG_q^\ast$ is $\cR$-periodic and its Bloch
decomposition is given by
\begin{equation}\label{eq:fibers_eiqx}
(G_qAG_q^\ast)_k = A_{k-q}, \quad \mbox{ for a.a. } k \in \R^2.
\end{equation}


\subsection{The Bloch theorem}
Before attacking the well-posedness of the current
$j_\eps(t) = \underline {\rm Tr}(J
\gamma_\eps(t))$ for $\eps, t \neq 0$, we first study an easier special case.
\begin{proposition}[Bloch theorem]\label{prop:null_current} The current satisfies
\begin{align*}
& j_0(t) = 0, \quad \forall t \ge 0 \;\;\; \mbox{(no current in the absence of external field)}, \\
& j_\eps(0)=0, \quad \forall \eps \ge 0 \;\;\;  \mbox{(continuity of the current at $t=0$).}
\end{align*}
\end{proposition}
\begin{proof}
We have
\[
j_0(t)=j_\eps(0)=\frac1{4\pi^2}\int_\mathcal B
\Tr(\partial_\alpha H_kP_k)\,dk.
\]
Since $V$ is real, $H_{-k}=\overline{H_k}$ and $P_{-k}=\overline{P_k}$,
where the bar denotes complex conjugation. Moreover,
$\partial_\alpha H_{-k}=-\overline{\partial_\alpha H_k}$.
Thus the real-valued, $\cR^*$-periodic function
$k\mapsto\Tr(\partial_\alpha H_kP_k)$ is odd, and its integral vanishes.
\end{proof}


\subsection{Definition of the current}
For $\eps>0$, the operator $H_\eps=H-\eps x_\beta$ is not
$\cR$-periodic, and we would naively expect that the density matrix
$\gamma_\eps(t)$ is not either. Yet, this operator is in fact
$\cR$-periodic. Physically, this is due to the fact that although
$-\eps x\cdot e_\beta$ is not periodic, its negative gradient is constant,
hence periodic. The proof relies on the standard gauge transform
\begin{align}\label{eq:gauge_change}
\widetilde\psi(x,t)=G_{-\eps t e_\beta}\psi(x,t)
=e^{-i\eps t x_\beta}\psi(x,t),
\end{align}
where $G_q$ was defined in~\eqref{eq:defGq}, and the gauge-transformed operators
\begin{align}\label{gauge_u}
U_\eps(t,t')=G_{-\eps t e_\beta}
 e^{-i(t-t')H_\eps}G_{-\eps t'e_\beta}^*,
\end{align}
and
\begin{equation}\label{gauge}
\widetilde\gamma_\eps(t)=G_{-\eps t e_\beta}\gamma_\eps(t)G_{-\eps t e_\beta}^*
=U_\eps(t)\gamma_0U_\eps(t)^*,
\end{equation}
where $U_\eps(t)=U_\eps(t,0)$. Through this change of gauge, the dynamics
induced by $H_\eps$ is equivalent to that induced by the time-dependent
periodic Hamiltonian
\begin{align}\label{widetilde_H}
\widetilde H_\eps(t)=G_{-\eps t e_\beta}H_\eps G_{-\eps t e_\beta}^*
+\eps x_\beta=\frac12(p+\eps t e_\beta)^2+V.
\end{align}
Its Bloch fibers are
\begin{align}\label{widetilde_H_k}
\widetilde H_{\eps,k}(t)=\frac12(p+k+\eps t e_\beta)^2+V=H_{k+\eps t e_\beta}.
\end{align}
We sum up these arguments in the proposition below, together with elements
that we shall use for defining the current.

\begin{proposition}\label{prop:setting}
Let $V\in H^1_{\rm per}(\Omega;\R)$.
\begin{enumerate}
\item For every $\eps\in\R$, $H_\eps$ is essentially self-adjoint on
$C_c^\infty(\R^2;\C)$ and therefore generates a unitary group on $L^2(\R^2;\C)$.
\item The operator $\widetilde H_\eps(t)$ is self-adjoint with domain
$H^2(\R^2;\C)$ and $\cR$-periodic. The propagator $U_\eps(t,t')$ is
$\cR$-periodic, with fibers $U_{\eps,k}(t,t')$ solving
\begin{equation}\label{eq:propagator2}
i\partial_tU_{\eps,k}(t,t')=H_{k+\eps te_\beta}U_{\eps,k}(t,t'),
\qquad U_{\eps,k}(t',t')=\mathrm{Id}.
\end{equation}
\item For $t\ge0$, $J\gamma_\eps(t)\in\Schatten_{1,\rm per}$, and
\begin{align}\label{eq:current-k}
j_\eps(t)&=\frac1{4\pi^2}\int_\cB
\Tr\left(\partial_\alpha H_{k+\eps te_\beta}\widetilde\gamma_{\eps,k}(t)\right)\,dk\\
&=\frac1{4\pi^2}\int_\cB\Tr\left(\partial_\alpha H_{k+\eps te_\beta}
U_{\eps,k}(t)P_kU_{\eps,k}(t)^*\right)\,dk.\nonumber
\end{align}
\end{enumerate}
\end{proposition}
The results of Proposition~\ref{prop:setting} are proved in
Section~\ref{sec:preliminaries} for completeness. The situation can be
summed up in the commutative diagrams of Figure~\ref{fig:diagram}.
\begin{figure}[h!]
\centering
\begin{tikzcd}[column sep=large,row sep=large]
\gamma_0\arrow{r}{e^{-itH_\eps}}\arrow{d}{\mathrm{Id}}&
\gamma_\eps(t)\arrow{d}{G_{-\eps te_\beta}}\\
\gamma_0\arrow{r}{U_\eps(t)}&\widetilde\gamma_\eps(t)
\end{tikzcd}
\medskip

\begin{tikzcd}[column sep=large,row sep=large]
\{P_k\}_k\arrow{r}{}\arrow{d}{\mathrm{Id}}&
\{\gamma_{\eps,k}(t)\}_k\arrow{d}{k\mapsto k+\eps te_\beta}\\
\{P_k\}_k\arrow{r}{\{U_{\eps,k}(t)\}_k}&
\{\widetilde\gamma_{\eps,k}(t)\}_k
\end{tikzcd}
\caption{Relations between the density matrices and their Bloch fibers.
An arrow labelled by a unitary $U$ denotes conjugation by $U$.
The right lower arrow means
$\widetilde\gamma_{\eps,k}(t)=\gamma_{\eps,k+\eps te_\beta}(t)$.}
\label{fig:diagram}
\end{figure}


\subsection{Semimetals}\label{sec:semimetal-assumptions}
The position of the Fermi level in the band diagram is key to determining
the electronic properties of the medium. We assume the following.
\begin{assumption}[semimetal]\label{ass:semimetal}
There are $N,M\ge1$ and distinct points $k_1,\ldots,k_M$ modulo $\cR^*$ such that
\[
\lambda_{N,k}<\mu_{\rm F}<\lambda_{N+1,k}\qquad
(k\notin\{k_1,\ldots,k_M\}\bmod\cR^*).
\]
For each $i=1,\ldots,M$, the two bands meet at an isolated conical crossing:
\begin{align}
\lambda_{N-1,k_i}<\lambda_{N,k_i}=\mu_{\rm F}
=\lambda_{N+1,k_i}<\lambda_{N+2,k_i},\label{hyp:semimetals_1}\\
\lambda_{N,k}=\mu_{\rm F}-v_i|k-k_i|+O(|k-k_i|^2),\label{hyp:semimetals_2}\\
\lambda_{N+1,k}=\mu_{\rm F}+v_i|k-k_i|+O(|k-k_i|^2),\label{hyp:semimetals_3}
\end{align}
where $v_i>0$.
\end{assumption}
The Fermi surface is the finite set
\begin{equation}\label{eq:Fermi_surface}
\mathcal S=\{k_1,\ldots,k_M\}\bmod\cR^*.
\end{equation}
For $k\notin\mathcal S$, $P_k=P_{N,k}$ has rank $N$. Its values at the
crossings do not affect any Brillouin-zone integral. Since $V$ is real,
$H_{-k}=\overline{H_k}$: the system has time-reversal symmetry.
The passage to linear response is justified in
Proposition~\ref{prop:response}.


\subsection{Conductivity}
In the following, $O(f(\eps,t))$ denotes a quantity bounded by $Cf(\eps,t)$,
where $C$ may depend on the potential, the Fermi level and the fixed
directions, but not on $t$ or $\eps$. A subscript $T$ allows the constant
to depend on a fixed time horizon $T$. We write $o_T(\eps)$ for a
remainder whose quotient by $\eps$ tends uniformly to zero on that interval.
\begin{theorem}\label{theo:semi-metal}
Under Assumption~\ref{ass:semimetal}, the iterated limit exists and
\[
\sigma=\lim_{T\to\infty}\frac1T\int_0^T
\left(\lim_{\eps\to0^+}\frac{j_\eps(t)}{\eps}\right)\,dt
=\frac{M}{16}\,e_\alpha\cdot e_\beta.
\]
\end{theorem}
The value of the conductivity is universal: it does not depend on the
characteristics of the Hamiltonian, but only on the number of conical
crossings. The conductivity tensor is isotropic and each conical
intersection contributes $1/16$.


\section{Bloch decomposition and regularity of the current}\label{sec:preliminaries}
In this section, we prove Proposition~\ref{prop:setting}. We first
point out an alternative interpretation of the gauge change.
Formally, $\gamma_\eps(t)$ satisfies
\[
i\partial_t\gamma_\eps=[H,\gamma_\eps]-\eps[x_\beta,\gamma_\eps].
\]
The commutator $[x_\beta,\gamma_\eps]$ is $\cR$-periodic, with fibers
$i\partial_\beta\gamma_{\eps,k}$. Therefore the fibers satisfy
\[
i\partial_t\gamma_{\eps,k}+i\eps\partial_\beta\gamma_{\eps,k}
=[H_k,\gamma_{\eps,k}]=L_{H_k}\gamma_{\eps,k},
\]
where $L_{H_k}=[H_k,\cdot]$ is the Liouvillian associated with $H_k$
(see Section~\ref{sec:louvillian_intro}). The left-hand side is a linear
advection equation, which suggests the method of characteristics: setting
\begin{equation}\label{eq:change_of_variable}
\widetilde\gamma_{\eps,k}(t)=\gamma_{\eps,k+\eps te_\beta}(t),
\end{equation}
we obtain
\[
i\partial_t\widetilde\gamma_{\eps,k}(t)
=[H_{k+\eps te_\beta},\widetilde\gamma_{\eps,k}(t)].
\]
The gauge transform $G_{-\eps te_\beta}$, equivalent to the change of
variable~\eqref{eq:change_of_variable}, makes these remarks rigorous.
We now prove Proposition~\ref{prop:setting}.


{\bf Proof of the first assertion.}
The essential self-adjointness of $H_\eps$ follows from the
Faris--Lavine theorem~\cite[Theorem~X.38]{ReedSimon2}, applied to
$2H_\eps=-\Delta+2V-2\eps x_\beta$. Let
$\mathcal C=C_c^\infty(\R^2;\C)$. We have
\[
-2\eps x\cdot e_\beta\ge-|x|^2-\eps^2|e_\beta|^2.
\]
In two dimensions, $V\in L^2_{\rm per}$ is infinitesimally bounded with
respect to $-\Delta$~\cite[Theorem~XIII.96]{ReedSimon4}, hence also with
respect to the harmonic oscillator $-\Delta+2|x|^2$.
The linear potential $-2\eps x_\beta$ is likewise
infinitesimally bounded with respect to this oscillator. The Kato--Rellich
theorem therefore gives the essential self-adjointness of
$-\Delta+2V-2\eps x_\beta+2|x|^2$ on the core $\mathcal C$.
Also, the infinitesimal form bound for $V$ gives that $-a\Delta+2V$
is bounded below for every $a>0$. Since $\mathcal C$ is invariant
under multiplication by coordinates and differentiation, all the
hypotheses of the Faris--Lavine theorem are satisfied.
Thus $H_\eps$ is essentially self-adjoint on $\mathcal C$, and its
propagator is the unitary group
$(e^{-itH_\eps})_{t\in\R}$.

\medskip

{\bf Proof of the second assertion.}
The self-adjointness of
\[
\widetilde H_\eps(t)=\frac12(p+\eps te_\beta)^2+V
\]
on $H^2(\R^2;\C)$ follows from the Kato--Rellich theorem.
To show the well-posedness of the dynamics, it suffices to study its
fibers $H_{k+\eps te_\beta}$. For $t\in[0,T]$, write
$H_{k+\eps te_\beta}=H_0+H_1(t)$, with $H_0=-\frac12\Delta$ and
\[
H_1(t)=(k+\eps te_\beta)\cdot p+\frac12|k+\eps te_\beta|^2+V.
\]
The perturbation $H_1(t)$ is infinitesimally $H_0$-bounded, uniformly
when $k$, $\eps$ and $t$ range over bounded sets. Hence the Hamiltonians
have common domain $H^2_{\rm per}$ and graph norms uniformly equivalent
to the $H^2_{\rm per}$-norm. Their dependence on $t$ is continuously
differentiable as operators from $H^2_{\rm per}$ to $L^2_{\rm per}$;
indeed,
\[
H_{k+\eps te_\beta}-H_{k+\eps se_\beta}
=\eps(t-s)\partial_\beta H_k
+\frac{\eps^2}{2}(t^2-s^2)|e_\beta|^2.
\]
After adding a fixed scalar to the Hamiltonians to make them invertible,
these facts verify the resolvent and difference-quotient hypotheses of
the common-domain propagator theorem~\cite[Theorem~X.70]{ReedSimon2}.
It gives a unique unitary propagator $U_{\eps,k}(t,s)$ whose trajectories
from $H^2_{\rm per}$ belong to
$C([0,T];H^2_{\rm per})\cap C^1([0,T];L^2_{\rm per})$.
The graph-norm estimate in that theorem also gives
\begin{equation}\label{eq:propagator_H2}
\|U_{\eps,k}(t,s)\|_{\mathcal L(H^2_{\rm per})}\le C_T,
\qquad k\in\mathcal B,\quad 0\le s\le t\le T,\quad 0\le\eps\le1.
\end{equation}

The same theorem applies to $\widetilde H_\eps(t)$ on
$L^2(\R^2;\C)$, and its propagator is the direct integral of the fiber
propagators. Undoing the gauge change~\eqref{eq:gauge_change} in the
Schr\"odinger equation, tested against $C_c^\infty$ functions, identifies
it with~\eqref{gauge_u} by uniqueness for the self-adjoint generator
$H_\eps$. This proves the second assertion.


\noindent
{\bf Proof of the third assertion.} We first compute the fibers of
$\gamma_\eps(t)$. Using~\eqref{eq:fibers_eiqx}, we have
\begin{align}
\gamma_{\eps,k}(t)
&=\left(G_{-\eps te_\beta}^*U_\eps(t)\gamma_0U_\eps(t)^*G_{-\eps te_\beta}\right)_k\\
&=U_{\eps,k-\eps te_\beta}(t)P_{k-\eps te_\beta}
 U_{\eps,k-\eps te_\beta}(t)^*\label{eq:gebkt}\\
&=\sum_{n:\,\lambda_{n,k-\eps te_\beta}\le\mu_{\rm F}}
 |U_{\eps,k-\eps te_\beta}(t)u_{n,k-\eps te_\beta}\rangle
 \langle U_{\eps,k-\eps te_\beta}(t)u_{n,k-\eps te_\beta}|,\nonumber
\end{align}
for almost every $k$. Since $u_{n,k}\in H^2_{\rm per}$, we deduce that
\[
(J\gamma_\eps(t))_k=\partial_\alpha H_k\gamma_{\eps,k}(t)
\]
is trace-class and finite-rank, with a uniform bound for $k\in\cB$.
The current is therefore well-defined. Since
$k\mapsto\Tr(\partial_\alpha H_k\gamma_{\eps,k}(t))$ is $\cR^*$-periodic,
\begin{align*}
j_\eps(t)
&=\frac1{4\pi^2}\int_\cB\Tr(\partial_\alpha H_k\gamma_{\eps,k}(t))\,dk\\
&=\frac1{4\pi^2}\int_\cB\Tr(\partial_\alpha H_{k+\eps te_\beta}
\gamma_{\eps,k+\eps te_\beta}(t))\,dk,
\end{align*}
which gives~\eqref{eq:current-k}.


\section{Linear response}\label{sec:linear-response}

\subsection{The Liouvillian and its partial inverse}
\label{sec:louvillian_intro}
In order to formulate and interpret our results, it is convenient to make use of the Liouvillian formalism. For each $k\notin\mathcal S$, the Liouvillian
$L_k$ associated with $H_k$ is the commutator
\begin{equation}\label{eq:Liouvillian}
L_k A=[H_k,A].
\end{equation}
The operator $L_k$ is to density matrices what the Hamiltonian $H_k$ is
to pure states: its evolution is
\begin{equation}\label{eq:Liouvillian2}
e^{-itL_k}A=e^{-itH_k}Ae^{itH_k}.
\end{equation}
More precisely, $L_k$ is the self-adjoint generator of this
strongly-continuous unitary group on the Hilbert--Schmidt operators.
Its action is easily understood in the eigenbasis of $H_k$:
\[
L_k|u_{n,k}\rangle\langle u_{m,k}|=
(\lambda_{n,k}-\lambda_{m,k})|u_{n,k}\rangle\langle u_{m,k}|.
\]
The operator $L_k$ is not invertible. However, it is invertible
on the block off-diagonal operators, i.e. operators $A$ such that
$P_kAP_k=(1-P_k)A(1-P_k)=0$. Its partial inverse $L_k^+$ is given by
\begin{align}\label{eq:explicit_inverse_liouvillian}
L_k^+A=\sum_{n\le N<m}
\frac{\langle u_{m,k},Au_{n,k}\rangle
 |u_{m,k}\rangle\langle u_{n,k}|
-\langle u_{n,k},Au_{m,k}\rangle
 |u_{n,k}\rangle\langle u_{m,k}|}
{\lambda_{m,k}-\lambda_{n,k}}.
\end{align}
The gap and the finite number of occupied states ensure convergence in
Hilbert--Schmidt norm for bounded $A$. We extend $L_k^+$ to all bounded
operators by setting it to zero on the diagonal blocks. In particular,
$L_kL_k^+A=A$ for off-diagonal $A$.

\subsection{Spectral derivatives}
For $k\notin\mathcal S$, the contour formula
\begin{equation}\label{eq:int_rep}
P_k=\frac{1}{2\pi i}\oint_{\mathcal C}(z-H_k)^{-1}\,dz,
\end{equation}
where $\mathcal C$ encloses the occupied eigenvalues, gives a locally
analytic projector. Differentiating $[H_k,P_k]=0$, we get
\[
L_k\partial_\beta P_k=[P_k,\partial_\beta H_k].
\]
Both sides are off-diagonal (for $\partial_\beta P_k$, differentiate
$P_k^2=P_k$). Hence
\begin{equation}\label{eq:Pprime}
\partial_\beta P_k=L_k^+[P_k,\partial_\beta H_k].
\end{equation}

\subsection{Linear response}
\label{subsec:lr}

We now aim at obtaining an expansion of the current to first order in
$\eps$ for a given $t$. This is a classical computation in response
theory, sometimes known as the Kubo formula. We only need
differentiability at $\eps=0$ on bounded time intervals; the estimates
below do not require a lower bound on the gap.

\begin{proposition}\label{prop:response}
For every $T>0$, uniformly for $k\notin\mathcal S$ and $0\le t\le T$,
\begin{align}\label{eq:response_expansion}
\Tr\left(\partial_\alpha H_{k+\eps te_\beta}
U_{\eps,k}(t)P_kU_{\eps,k}(t)^*\right)
={}&\partial_\alpha\Tr(H_kP_k)
+\eps t\partial_\alpha\partial_\beta\Tr(H_kP_k)\nonumber\\
&-i\eps\Tr\left(\partial_\alpha H_k(e^{-itL_k}-1)
L_k^+\partial_\beta P_k\right)+o_T(\eps).
\end{align}
The current may be differentiated at $\eps=0$ under the Bloch integral.
The limit defining $j^{\rm LR}(t)$ exists and may also be passed under
any integral over a bounded time interval.
\end{proposition}

\begin{proof}
We first justify differentiation with the unbounded current observable.
Fix $T>0$ and $0\le\eps\le1$. By~\eqref{eq:propagator_H2}, the
propagators are uniformly bounded on $H^2_{\rm per}$ for
$0\le s\le t\le T$ and $k\in\mathcal B$.
The occupied unit vectors also have uniformly bounded $H^2_{\rm per}$
norms, since their energies are bounded.

For such a vector $u$, set
$\psi_\eps(t)=U_{\eps,k}(t)u$ and $\psi_0(t)=e^{-itH_k}u$.
Since
\[
H_{k+\eps se_\beta}-H_k
=\eps s\partial_\beta H_k+\frac{\eps^2s^2|e_\beta|^2}{2},
\]
Duhamel's formula gives
\begin{equation}\label{eq:duhamel_current}
\psi_\eps(t)-\psi_0(t)
=-i\int_0^t e^{-i(t-s)H_k}
\left(\eps s\partial_\beta H_k+\frac{\eps^2s^2|e_\beta|^2}{2}\right)
\psi_\eps(s)\,ds.
\end{equation}
The difference is $O_T(\eps)$ in $L^2_{\rm per}$ and bounded in
$H^2_{\rm per}$, hence $O_T(\sqrt\eps)$ in $H^1_{\rm per}$ by
interpolation. Dividing~\eqref{eq:duhamel_current} by $\eps$ therefore
gives, in $L^2_{\rm per}$,
\[
\frac{\psi_\eps(t)-\psi_0(t)}{\eps}
\longrightarrow-i\int_0^t s\,e^{-i(t-s)H_k}
\partial_\beta H_ke^{-isH_k}u\,ds.
\]
All these bounds and the convergence are uniform in $k$, in occupied
unit vectors $u$, and in $0\le t\le T$.

Put $w=\psi_\eps-\psi_0$. The identity
\[
\langle\psi_\eps,\partial_\alpha H_k\psi_\eps\rangle
-\langle\psi_0,\partial_\alpha H_k\psi_0\rangle
=2\operatorname{Re}\langle w,\partial_\alpha H_k\psi_0\rangle
+\langle w,\partial_\alpha H_kw\rangle
\]
has last term $O_T(\eps^{3/2})$. Thus the current expectations are
differentiable, with a uniform $o_T(\eps)$ remainder. Summing over the
$N$ occupied vectors gives the density-matrix derivative
\begin{equation}\label{eq:density_derivative}
\left.\frac{d}{d\eps}\widetilde\gamma_{\eps,k}(t)\right|_{\eps=0}
=-i\int_0^t s\,e^{-i(t-s)L_k}[\partial_\beta H_k,P_k]\,ds.
\end{equation}
This formula also holds when paired with $\partial_\alpha H_k$:
$\partial_\beta H_ku\in H^1_{\rm per}$, and the free propagator
preserves this form domain.

We now use~\eqref{eq:Pprime}. Since
$[\partial_\beta H_k,P_k]=-L_k\partial_\beta P_k$ and
$\partial_\beta P_k$ is off-diagonal, integration by parts gives
\[
-i\int_0^t s\,e^{-i(t-s)L_k}[\partial_\beta H_k,P_k]\,ds
=t\partial_\beta P_k-i(e^{-itL_k}-1)L_k^+\partial_\beta P_k.
\]
By the Hellmann--Feynman theorem,
\begin{align*}
\partial_\alpha\Tr(H_kP_k)&=\Tr(\partial_\alpha H_kP_k),\\
\partial_\alpha\partial_\beta\Tr(H_kP_k)
&=N e_\alpha\cdot e_\beta
+\Tr(\partial_\alpha H_k\partial_\beta P_k).
\end{align*}
Together with
$\partial_\alpha H_{k+\eps te_\beta}
=\partial_\alpha H_k+\eps t e_\alpha\cdot e_\beta$,
these identities prove~\eqref{eq:response_expansion}.
The uniform bounds justify the stated limit interchanges, and
Proposition~\ref{prop:null_current} removes the equilibrium current.
\end{proof}


\section{Proof of the conductivity formula}\label{sec:semimet}

We prove here Theorem~\ref{theo:semi-metal}. We first take the limit
$\eps\to0^+$ in Proposition~\ref{prop:response}, and cancel the term
proportional to $t$. After averaging in time, only the two crossing
bands near each Dirac point contribute to the limit. We construct
their reduced Hamiltonian and compare it with the Dirac Hamiltonian,
for which we can explicitly compute the limiting current.

\subsection{Cancellation of the term linear in time}
\label{sec:band_cancellation}

For $n\le N<m$, write
\begin{align}\label{eq:transition_data}
\Delta_{nm,k}&=\lambda_{m,k}-\lambda_{n,k},\nonumber\\
z_{nm,k}&=\langle u_{n,k},\partial_\beta H_ku_{m,k}\rangle
\langle u_{m,k},\partial_\alpha H_ku_{n,k}\rangle.
\end{align}
The finite number of occupied states, the uniform bounds on
$\partial_\alpha H_ku_{n,k}$ and $\partial_\beta H_ku_{n,k}$,
and Cauchy--Schwarz give
\begin{equation}\label{eq:transition_bound}
\sum_{n\le N<m}|z_{nm,k}|\le C.
\end{equation}
This controls the infinitely many unoccupied bands. Sums over complete
eigenspaces are independent of the choice of basis at internal band crossings.

Let $F(k)=\Tr(H_kP_k)$. The spectral derivative~\eqref{eq:Pprime}
and the eigenbasis representation~\eqref{eq:explicit_inverse_liouvillian}
give
\[
\partial_\alpha\partial_\beta F(k)
=N e_\alpha\cdot e_\beta
-2\sum_{n\le N<m}\frac{\operatorname{Re}z_{nm,k}}{\Delta_{nm,k}}.
\]
The gap is bounded below away from the crossings and is comparable
to $|k-k_i|$ near $k_i$. Thus this derivative is integrable.
Applying Stokes formula on the momentum torus with disks of radius $r$
removed gives
\begin{equation}\label{eq:filled_band_cancellation}
\int_{\mathcal B}\partial_\alpha\partial_\beta F(k)\,dk=0.
\end{equation}
Indeed, $\partial_\alpha F(k)=\Tr(\partial_\alpha H_kP_k)$ is bounded,
so the extra boundary integrals are $O(r)$ and vanish as $r\to0$.

For the remaining term of~\eqref{eq:response_expansion}, we use
\begin{align*}
&\Tr\left(\partial_\alpha H_k(e^{-itL_k}-1)L_k^+\partial_\beta P_k\right)\\
&\qquad=\sum_{n\le N<m}(e^{it\Delta_{nm,k}}-1)
\frac{z_{nm,k}}{\Delta_{nm,k}^2}-{\rm c.c.}
\end{align*}
Time reversal leaves the gaps unchanged and conjugates the products
$z_{nm,k}$. Their imaginary part therefore integrates to zero.
Using~\eqref{eq:filled_band_cancellation}, we obtain
\begin{equation}\label{eq:linear_response_current}
j^{\rm LR}(t)=\frac1{2\pi^2}\int_{\mathcal B}
\sum_{n\le N<m}\operatorname{Re}z_{nm,k}
\frac{\sin(t\Delta_{nm,k})}{\Delta_{nm,k}^2}\,dk.
\end{equation}
The preceding decompositions are legitimate:
$|\sin(t\Delta)|/\Delta^2\le t/\Delta$ and
$|1-\cos(t\Delta)|/\Delta^2\le t^2/2$, so every term is integrable
by~\eqref{eq:transition_bound} and the conical gap bound.

\subsection{Time averaging and reduction to two bands}
\label{sec:dirac_reduc}

We can now average in time. By Fubini's theorem,
\begin{equation}\label{eq:averaged_response}
\frac1T\int_0^Tj^{\rm LR}(t)\,dt
=\frac1{2\pi^2}\int_{\mathcal B}
\sum_{n\le N<m}\operatorname{Re}z_{nm,k}K_T(\Delta_{nm,k})\,dk,
\qquad K_T(\Delta)=\frac{1-\cos(T\Delta)}{T\Delta^3}.
\end{equation}
This formula makes the large-time behavior explicit: a gap bounded
away from zero contributes $O(T^{-1})$, while near a Dirac point the
relevant momenta have size $T^{-1}$.

\medskip

We introduce a small parameter $\delta>0$ controlling the size of the
neighborhood of the Dirac points, which is independent of $T$.
We decompose $\mathcal B$ as the disjoint union
\[
\mathcal B=\mathcal B_{\rm out}^{\delta}
\;\cup\;(\cup_{i=1}^M\mathcal B_i^\delta),
\qquad \mathcal B_i^\delta=B(k_i,\delta),
\]
with all sets understood modulo the reciprocal lattice. The radius is
chosen small enough that the balls are disjoint and the two crossing
bands are separated from the other bands in each ball.
We keep $\delta$ fixed throughout the proof.

On $\mathcal B_{\rm out}^\delta$, the energy differences
$\Delta_{nm,k}$ are bounded away from zero for each fixed $\delta$.
Since $|K_T(\Delta)|\le2/(T\Delta^3)$, the bound
\eqref{eq:transition_bound} shows that the contribution of this region
to~\eqref{eq:averaged_response} is $O_\delta(T^{-1})$.
Here $O_\delta$ permits a constant depending on the fixed $\delta$.

When $n\ne N$ or $m\ne N+1$, the denominators in
\eqref{eq:averaged_response} are bounded from below in the balls,
independently of $\delta$. These transitions also contribute
$O(T^{-1})$. Putting the two regions together, we obtain
\begin{equation}\label{eq:two_band_average}
\frac1T\int_0^Tj^{\rm LR}(t)\,dt
=\frac1{2\pi^2}\sum_{i=1}^M I_i^R(\delta,T)+O_\delta(T^{-1}),
\end{equation}
where
\begin{equation}\label{eq:IRdt}
I_i^R(\delta,T)=\int_{\mathcal B_i^\delta}
\operatorname{Re}z_{N,N+1,k}
\frac{1-\cos(T\Delta_{N,N+1,k})}{T\Delta_{N,N+1,k}^3}\,dk.
\end{equation}
At this stage of the proof, only two modes are involved in the formula
giving the current, namely the two modes that cross at the Fermi level.
Everything happens as for a two-band model that we now study.


\subsection{The local two-band model}
We now are interested in the computation of $I_i^R(\delta,T)$.
In the following, we drop the index $i$ and assume without loss of
generality that $k_i=0$. We write $v=v_i$.

Hypothesis \eqref{hyp:semimetals_1} implies that, for $k$ small
enough, the Bloch Hamiltonian $H_k$ has exactly two eigenvalues
(counting multiplicities) close to $\mu_{\rm F}$. Consider an
arbitrary orthonormal basis $(v_0,w_0)$ of
$\Ran(P_{N+1,0} - P_{N-1,0})$. For all $k$ small
enough, we can construct an orthonormal basis $(v_k,w_k)$ of
$\Ran(P_{N+1,k} - P_{N-1,k})$ by L\"owdin
orthonormalization of
$((P_{N+1,k} - P_{N-1,k})v_0,(P_{N+1,k} -
P_{N-1,k})w_0)$, and set
$$
H_k^{R} = [v_{k} | w_{k}]^{*}
H_k[v_{k} | w_{k}] = \left( \begin{array}{cc} \langle v_k, H_k v_k \rangle &  \langle v_k, H_k w_k \rangle \\
\langle w_k, H_k v_k \rangle &  \langle w_k, H_k w_k \rangle \end{array} \right).
$$
It follows that the reduced Hamiltonian
\begin{align*}
  H_k\Big|_{\Ran(P_{N+1,k} - P_{N-1,k}) }
\end{align*}
is equivalent through a unitary transform that depends analytically on $k$
to the reduced $2 \times 2$ Hamiltonian
\begin{align}
  \label{eq:Hmatk}
  H_k^{R} = \sum_{p=0}^3 b^{p}(k) \sigma_{p},
\end{align}
where $$\sigma_{0} = {\rm Id}_{\C^{2}}, \;\; \sigma_{1}=\left( \begin{array}{cc} 0 & 1 \\ 1 & 0 \end{array} \right), \;\;\sigma_{2}=\left( \begin{array}{cc} 0 & -i \\ i & 0 \end{array} \right), \;\;\sigma_{3}=\left( \begin{array}{cc} 1 & 0 \\ 0 & -1 \end{array} \right)$$
are the Pauli matrices, and $(b^{p})_{p=0,1,2,3}$ are real-valued
analytic functions of $k$ in a neighborhood of $0$.

The matrix $H_k^{R}$ has eigenvalues
$$  \lambda_{\pm}(k) = b^{0}(k) \pm \sqrt{\sum_{p=1}^{3}b^{p}(k)^{2}}.
$$
It follows that
\begin{align*}
  b^{0}(k) &= \mu_{F} + O(|k|^{2}), \quad
  b^{p}(k) = v\langle q^{p}, k\rangle + O(|k|^{2}), \quad p=1,2,3,
\end{align*}
where the $(q^{p})_{p=1,2,3}$ are the rows of a $3 \times 2$
matrix $Q$ with orthonormal columns, so that
\begin{align}\label{def:HR}
  H_k^{R} = \mu_{F} + v (Q k) \cdot \sigma + O(|k|^{2}).
\end{align}
Let $R \in {\rm SO}(3)$ be a rotation matrix that maps
${\rm Ran}(Q)$ to ${\rm Span}(e^0_{1},e^0_{2})$, where $(e^0_{1},e^0_{2},e^0_{3})$ is the canonical basis of $\R^3$. Write $R$ as a rotation through an angle $\theta$ about a unit vector $n$.
The unitary matrix
\[
U=\cos(\theta/2)\mathrm{Id}_2-i\sin(\theta/2)n\cdot\sigma
\]
satisfies, by the Pauli matrix identities,
\begin{align*}
  R_{pq} = \frac 1 2 \Tr(\sigma_{p} U \sigma_{q} U^{*}).
\end{align*}
It follows that
\begin{align*}
  \Tr(\sigma_{3} U H_k^{R} U^{*}) &= v\sum_{q=1}^{3} (Q k)_{q} \Tr(\sigma_{3} U \sigma_{q} U^{*}) + O(|k|^{2})\\
  &= 2 v\langle e^0_3,RQk\rangle+O(|k|^2)\\
  &= O(|k|^{2}).
\end{align*}
Up to a unitary transform, we can therefore assume $Q$ to be
a $2\times 2$ matrix in~\eqref{def:HR}.

\subsection{Reduction to the Dirac Hamiltonian}
\paragraph{Reduction to $H_k^{R}$}
For $k \neq 0$, let $\lambda^{R}_{\pm,k}$ be the larger and smaller
eigenvalues of $H_k^{R}$ respectively, and $u^{R}_{\pm,k}$
associated orthonormal eigenvectors in $\C^2$. We have
$\lambda_{-,k}^{R} = \lambda_{N,k}, \lambda_{+,k}^{R} =
\lambda_{N+1,k}$, and
\begin{align*}
  [v_{k} | w_{k}]
u_{+,k}^{R} = e^{i\theta_{+}(k)}u_{N + 1,k},\qquad[v_{k} | w_{k}]
u_{-,k}^{R} = e^{i\theta_{-}(k)}u_{N,k}
\end{align*}
for some phases $\theta_{\pm}(k) \in \R$. We have
\begin{align*}
  \partial_\alpha H_k^{R}
  &= [v_{k} | w_{k}]^{*}
\partial_\alpha H_k[v_{k} | w_{k}]+ \partial_\alpha [v_{k} | w_{k}]^{*} [v_{k} | w_{k}] H_k^{R} + H_k^{R}  [v_{k} | w_{k}]^{*}\partial_\alpha [v_{k} | w_{k}]
\\
&= [v_{k} | w_{k}]^{*}
\partial_\alpha H_k[v_{k} | w_{k}] + O(|k|)
\end{align*}
where we have used for the first line that $H_k$ commutes with
$P_{N+1,k} - P_{N-1,k} = [v_{k} | w_{k}]
[v_{k} | w_{k}]^{*}$, and for the second that $H_k^{R} = \mu_{F} {\rm Id}_{2} + O(|k|)$ and
$\partial_\alpha \Big([v_{k} | w_{k}]^{*} [v_{k} | w_{k}] \Big) =
\partial_\alpha {\rm Id}_{2} = 0$. We therefore obtain
\begin{align*}
  \langle u^{R}_{+,k},\partial_\alpha H_k^{R}u^{R}_{-,k}\rangle&= e^{-i(\theta_{+}(k)-\theta_{-}(k))}\langle u_{N+1,k},\partial_\alpha H_ku_{N,k}\rangle + O(|k|)\\
  \langle u^{R}_{-,k}, \partial_\beta H_k^{R} u^{R}_{+,k} \rangle&= e^{+i(\theta_{+}(k)-\theta_{-}(k))}\langle u_{N,k},\partial_\beta H_ku_{N+1,k}\rangle + O(|k|).
\end{align*}
The phases cancel in $z_{N,N+1,k}$. Replacing the physical velocities
by the derivatives of $H_k^R$ therefore changes this product by $O(|k|)$.

\paragraph{Reduction to $H_k^D$}
By standard results of perturbation theory \cite[Chapter XII]{ReedSimon4}
applied to $H_k^R=\mu_F+v(Qk)\cdot\sigma+O(|k|^2)$, with gap
greater than a constant multiple of $|k|$,
\[
\lambda^R_{\pm,k}=\mu_F+\lambda^Q_{\pm,k}+O(|k|^2),
\qquad u^R_{\pm,k}=u^Q_{\pm,k}+O(|k|),
\]
where the superscript $Q$ refers to eigenvalues and appropriately
chosen orthonormal eigenvectors of the Hamiltonian
\[
H_k^Q=v(Qk)\cdot\sigma.
\]
Write $\Delta(k)=\Delta_{N,N+1,k}$, $z(k)=z_{N,N+1,k}$, and set
\[
z^Q(k)=\langle u^Q_{-,k},\partial_\beta H_k^Q u^Q_{+,k}\rangle
\langle u^Q_{+,k},\partial_\alpha H_k^Q u^Q_{-,k}\rangle.
\]
The preceding estimates give, uniformly in direction,
\[
\Delta(k)=2v|k|+O(|k|^2),\qquad z(k)=z^Q(k)+O(|k|).
\]
The function $z^Q$ depends only on the direction of $k$.

We now make the change of variables $k=q/T$ in~\eqref{eq:IRdt}:
\begin{equation}\label{eq:scaled_cone}
I^R(\delta,T)=\int_{\R^2}\1_{\{|q|<T\delta\}}
\operatorname{Re}z(q/T)
\frac{1-\cos(T\Delta(q/T))}{(T\Delta(q/T))^3}\,dq.
\end{equation}
For each $q\ne0$, $T\Delta(q/T)\to2v|q|$ and
$z(q/T)\to z^Q(q)$. Moreover, on the support of the integrand,
$c|q|\le T\Delta(q/T)\le C|q|$ and $|z(q/T)|\le C$.
Using $1-\cos s\le\min(s^2/2,2)$, its absolute value is bounded by
\[
C\min\{|q|^{-1},|q|^{-3}\},
\]
which is integrable on $\R^2$. Dominated convergence therefore gives
\begin{equation}\label{eq:linear_cone_limit}
\lim_{T\to\infty}I^R(\delta,T)=I^Q,
\qquad
I^Q=\int_{\R^2}\operatorname{Re}z^Q(q)
\frac{1-\cos(2v|q|)}{(2v|q|)^3}\,dq.
\end{equation}
This change of scale handles the nonlinear gap inside the oscillatory
factor directly.

We perform the change of variable $q'=Qq$ (recall that $Q$ is orthogonal)
and obtain
\[
I^Q=(Qe_\beta)^T I^D(Qe_\alpha),
\]
where the coefficients of the $2\times2$ matrix $I^D$ are
\[
I_{ij}^D=\int_{\R^2}
\operatorname{Re}\left(
\langle u^D_{-,k},\partial_{k_i}H_k^D u^D_{+,k}\rangle
\langle u^D_{+,k},\partial_{k_j}H_k^D u^D_{-,k}\rangle\right)
\frac{1-\cos(2v|k|)}{(2v|k|)^3}\,dk,
\]
and the superscript $D$ refers to the Dirac Hamiltonian
\[
H_k^D=vk\cdot\sigma.
\]

\paragraph{The Dirac Hamiltonian $H_k^D$}
We finish by computing $I^D$ explicitly. Let
$k=r(\cos\theta,\sin\theta)$. We have
\[
\lambda^D_{\pm,k}=\pm vr,\qquad
u^D_{+,k}=\frac1{\sqrt2}\begin{pmatrix}1\\e^{i\theta}\end{pmatrix},
\qquad
u^D_{-,k}=\frac1{\sqrt2}\begin{pmatrix}-e^{-i\theta}\\1\end{pmatrix}.
\]
By an explicit calculation, we obtain
\[
\int_0^{2\pi}\operatorname{Re}\left(
\langle u^D_{-,k},\partial_{k_i}H_k^D u^D_{+,k}\rangle
\langle u^D_{+,k},\partial_{k_j}H_k^D u^D_{-,k}\rangle\right)d\theta
=\pi v^2\delta_{ij}.
\]
It follows that
\begin{align*}
I_{ij}^D
&=\pi v^2\delta_{ij}\int_0^\infty
\frac{1-\cos(2vr)}{(2vr)^3}\,r\,dr\\
&=\frac\pi4\delta_{ij}\int_0^\infty\frac{1-\cos s}{s^2}\,ds
=\frac{\pi^2}{8}\delta_{ij}.
\end{align*}
For the last integral, the identity
$s^{-2}=\int_0^\infty a e^{-as}\,da$ and Tonelli's theorem give
\[
\int_0^\infty\frac{1-\cos s}{s^2}\,ds
=\int_0^\infty a\left(\frac1a-\frac{a}{a^2+1}\right)da
=\frac\pi2.
\]
Since $Q$ is orthogonal, each crossing contributes
$\pi^2(e_\alpha\cdot e_\beta)/8$ to the limit of $I_i^R$.
We finally get from~\eqref{eq:two_band_average}
\[
\lim_{T\to\infty}\frac1T\int_0^T j^{\rm LR}(t)\,dt
=\frac1{2\pi^2}\,\frac{M\pi^2}{8}\,e_\alpha\cdot e_\beta
=\frac{M}{16}\,e_\alpha\cdot e_\beta,
\]
which proves the result.


\bibliographystyle{plain}
\bibliography{refs}
\end{document}
